%% =============================================================================
%% section5.tex — Section 5: Application to NSECE Childcare Data
%% =============================================================================

\section{Application to NSECE Childcare Data}\label{sec:application}

We now apply the full analytical pipeline developed in \cref{sec:model} to
the 2019 NSECE data described in \cref{sec:data}, demonstrating each
methodological component in a realistic survey setting. The simulation
evidence in \cref{sec:simulation} established that the sandwich correction
substantially improves frequentist coverage and that the DER diagnostic reliably
identifies parameters affected by informative sampling; the present
application illustrates these tools in practice. The analysis proceeds in
five parts: model comparison via LOO-CV (\cref{sec:modelcomp}),
population-average inference with the two-track reporting strategy and DER
diagnostic (\cref{sec:popavg}), state-specific heterogeneity and
cross-margin independence (\cref{sec:statevar}), marginal effect
decomposition (\cref{sec:ame}), and a brief examination of policy
moderation (\cref{sec:policyresults}). Together, these subsections showcase
the three methodological contributions: the two-track reporting strategy
that resolves prior domination, the DER diagnostic that quantifies
parameter-specific design effects, and the marginal effect decomposition
that translates hurdle model coefficients into policy-relevant quantities.

%% ---------------------------------------------------------------------------
%% 5.1 Model Comparison and Selection
%% ---------------------------------------------------------------------------

\subsection{Model comparison and selection}\label{sec:modelcomp}

We fit the five nested models summarized in~\cref{tab:model-hierarchy}---from
the pooled M0 to the full M3b with policy moderators---using the estimation
pipeline described in~\cref{alg:estimation}, implemented in Stan
\citep{CarpenterEtAl2017}. Each model is estimated with 4
chains of 3,500 iterations (1,500 warmup + 2,000 sampling), yielding 8,000
post-warmup draws per parameter.

\paragraph{Convergence.}
All five models satisfy standard diagnostics~\citep{Vehtari2021}:
$\hat{R} < 1.01$ for every monitored parameter, minimum bulk effective
sample size exceeds 400, and zero divergent transitions were recorded.
\cref{tab:convergence-summary} summarizes these statistics; the
non-centered parameterization~\citep{BetancourtGirolami2015} successfully
eliminates the funnel geometry. Full parameter-by-parameter diagnostics
appear in Supplementary Material~E.

\begin{table}[t]
\centering
\caption{Convergence diagnostics across all models. $\hat{R}_{\max}$ is the
  maximum potential scale reduction factor across all monitored parameters.
  $\ESS_{\min}$ reports the minimum bulk and tail effective sample sizes.
  All models achieve the recommended thresholds ($\hat{R} < 1.01$,
  $\ESS > 400$, zero divergences).}\label{tab:convergence-summary}
\smallskip
\begin{tabular}{@{}lccccc@{}}
\toprule
Model & $\hat{R}_{\max}$ & $\ESS_{\min}^{\mathrm{bulk}}$
  & $\ESS_{\min}^{\mathrm{tail}}$ & Divergences \\
\midrule
M0: Pooled         & 1.001 & 3{,}390 & 2{,}679 & 0 \\
M1: Random int.\   & 1.004 &    815  & 1{,}415 & 0 \\
M2: Block-diag.\ SVC & 1.004 & 1{,}078 & 1{,}244 & 0 \\
M3a: Cross-margin  & 1.004 & 1{,}385 & 1{,}257 & 0 \\
M3b: Policy mod.\  & 1.004 & 1{,}245 & 1{,}463 & 0 \\
\bottomrule
\end{tabular}
\end{table}

\paragraph{LOO-CV comparison.}
We compare models using Pareto-smoothed importance-sampling leave-one-out
cross-validation~\citep[\LOO-CV;][]{VehtariEtAl2017}, where ``leave-one-out''
refers to individual providers, not entire states.
The pointwise log-likelihood for each provider~$i$ includes both hurdle
components: the extensive-margin (Bernoulli) log-likelihood always
contributes, and the intensive-margin (zero-truncated beta-binomial)
log-likelihood contributes only for participants ($z_i = 1$).
\cref{tab:loo} presents the expected log pointwise predictive density
($\widehat{\elpd}_{\LOO}$) for each model.
For the survey-weighted specification M3b-W, the LOO computation uses the
unweighted nominal log-likelihood, because the pseudo-log-likelihood is not
a proper predictive density; the comparisons thus rank models on their
nominal predictive ability.
Because LOO-CV operates at the provider level while many inferential
targets are state-level parameters, the LOO ranking should be interpreted
as a guide to observation-level fit rather than a direct assessment of
state-level inference quality.
The progression reveals
diminishing returns: adding state random intercepts (M0$\to$M1) yields a
decisive improvement ($\Delta\widehat{\elpd} = +317.8$, SE $= 25.9$,
exceeding 12 standard errors), and allowing state-varying coefficients
(M1$\to$M2) provides a further significant gain ($+23.6$ incrementally).
Introducing the cross-margin covariance (M2$\to$M3a) does not improve
predictive performance ($\Delta\widehat{\elpd} < 2$), consistent with the
near-zero empirical cross-margin correlation and with
\cref{prop:sigma12-nonid}: since $\bSigma_{12}$ receives no likelihood
information, its value lies in inference rather than prediction. Adding
policy moderators (M3a$\to$M3b) likewise produces no predictive gain; the
$\bGamma$ matrices redistribute the sources of state heterogeneity without
changing the provider-level predictive distribution.

\begin{table}[t]
\centering
\caption{Leave-one-out cross-validation comparison for models M0--M3b.
  $\widehat{\elpd}_{\LOO}$ is the expected log pointwise predictive density
  estimated via Pareto-smoothed importance sampling
  \citep{VehtariEtAl2017}; $\Delta\widehat{\elpd} = \widehat{\elpd}_{\mathrm{M0}} - \widehat{\elpd}_{\mathrm{M}_j}$ measures the
  predictive gain over M0 (more negative indicates better prediction); SE($\Delta$) is the standard
  error of the difference. A model improvement is flagged as significant
  when $|\Delta\widehat{\elpd}| > 2 \times \mathrm{SE}(\Delta)$.
  All Pareto-$k$ diagnostics are below 0.7.}\label{tab:loo}
\smallskip
\begin{tabular}{@{}lrcrrrc@{}}
\toprule
Model & Params & $\widehat{\elpd}_{\LOO}$ & SE &
  $\Delta\widehat{\elpd}$ & SE($\Delta$) & Sig.\ \\
\midrule
M0: Pooled          &  11 & $-20{,}945.0$ & 140.3 &
  ---   & ---  & ---    \\
M1: Random int.\    & 116 & $-20{,}627.2$ & 144.0 &
  $-317.8$ & 25.9 & Yes \\
M2: Block-diag.\ SVC & 551 & $-20{,}603.5$ & 144.4 &
  $-341.4$ & 27.6 & Yes \\
M3a: Cross-margin   & 576 & $-20{,}601.6$ & 144.4 &
  $-343.4$ & 27.6 & No  \\
M3b: Policy mod.\   & 616 & $-20{,}607.6$ & 144.6 &
  $-337.3$ & 28.5 & No  \\
\bottomrule
\end{tabular}
\end{table}

\paragraph{Model selection rationale.}
We adopt M3b as the reporting model despite its equivalent predictive
performance to M2 and M3a. The cross-margin covariance $\bSigma_{12}$
enables joint posterior statements about the reversal (e.g.,
$\Pr(\tilde{\alpha}_{\mathrm{pov},s} < 0$ and
$\tilde{\beta}_{\mathrm{pov},s} > 0 \mid \mathrm{data})$), and the policy
moderators provide the decomposition of state heterogeneity into
policy-explained and residual components (\cref{prop:poverty-expansion}).
The non-significant LOO-CV differences confirm that this additional
structure incurs no predictive cost. Posterior predictive checks confirm
that M3b reproduces both the observed zero rate (35.3\%) and the
distribution of IT shares among servers; details appear in Supplementary
Material~D.

%% ---------------------------------------------------------------------------
%% 5.2 Population-Average Effects and the Sandwich Correction
%% ---------------------------------------------------------------------------

\subsection{Population-average effects and the sandwich correction}\label{sec:popavg}

\cref{tab:fixed-effects} reports the posterior means, naive 95\% credible
intervals (from the raw MCMC output), sandwich-corrected 95\% Wald confidence
intervals, and design effect ratios (DER) for all 11 fixed-effect parameters.
Of the 11 parameters, 9 are statistically significant at the 95\% level
under the sandwich correction; the two exceptions are
$\alpha_{\mathrm{Hisp}}$ (extensive-margin Hispanic coefficient, Wald CI
$[-0.095, 0.106]$) and $\beta_{\mathrm{Black}}$ (intensive-margin Black
coefficient, Wald CI $[-0.052, 0.019]$).

\begin{table}[t]
\centering
\caption{Population-average fixed effect estimates from weighted M3b.
  ``Post.\ Mean'' is the posterior mean from the pseudo-posterior. ``Naive
  95\% CI'' is the 2.5th--97.5th percentile interval from the uncorrected
  MCMC posterior. ``Wald 95\% CI'' uses the sandwich-corrected variance
  (\cref{thm:cholesky}). ``DER'' is the design effect ratio
  (\cref{eq:der}). Asterisks indicate significance at the 5\% level
  under the Wald test. All covariates are standardized.}\label{tab:fixed-effects}
\smallskip
\small
\begin{adjustbox}{max width=\textwidth}
\begin{tabular}{@{}lrll r@{}}
\toprule
Parameter & Post.\ Mean & Naive 95\% CI & Wald 95\% CI & DER \\
\midrule
\multicolumn{5}{@{}l}{\textit{Extensive margin ($\balpha$)}} \\[2pt]
Intercept           & $0.764$  & $[-1.109,\; 2.628]$ & $[0.677,\; 0.851]^{*}$  & 2.04 \\
Poverty             & $-0.324$ & $[-2.205,\; 1.497]$ & $[-0.426,\; -0.221]^{*}$ & 2.41 \\
Urban               & $0.442$  & $[-1.387,\; 2.224]$ & $[0.375,\; 0.509]^{*}$  & 4.18 \\
Black               & $0.478$  & $[-1.459,\; 2.405]$ & $[0.392,\; 0.565]^{*}$  & 1.54 \\
Hispanic            & $0.006$  & $[-1.902,\; 1.903]$ & $[-0.095,\; 0.106]$     & 1.82 \\[4pt]
\multicolumn{5}{@{}l}{\textit{Intensive margin ($\bbeta$)}} \\[2pt]
Intercept           & $-0.242$ & $[-2.058,\; 1.587]$ & $[-0.272,\; -0.213]^{*}$ & 1.36 \\
Poverty             & $0.090$  & $[-1.713,\; 1.904]$ & $[0.053,\; 0.127]^{*}$  & 1.61 \\
Urban               & $-0.047$ & $[-1.858,\; 1.768]$ & $[-0.087,\; -0.007]^{*}$ & 3.18 \\
Black               & $-0.017$ & $[-1.851,\; 1.760]$ & $[-0.052,\; 0.019]$     & 1.38 \\
Hispanic            & $-0.097$ & $[-1.974,\; 1.716]$ & $[-0.130,\; -0.064]^{*}$ & 1.14 \\[4pt]
\multicolumn{5}{@{}l}{\textit{Overdispersion}} \\[2pt]
$\log\kappa$        & $1.919$  & $[1.866,\; 1.972]$ & $[1.861,\; 1.976]^{*}$   & 2.58 \\
\bottomrule
\end{tabular}
\end{adjustbox}
\end{table}

\paragraph{Impact of survey weighting.}
\cref{tab:m3b-comparison} places the weighted (M3b-W) and unweighted
(M3b) posterior means side by side.  Weighting shifts every
fixed-effect estimate, with a median displacement of 5.2 sandwich
standard errors, yet preserves all substantively important sign
patterns---most critically, the poverty reversal
($\alpha_{\mathrm{pov}} < 0$, $\beta_{\mathrm{pov}} > 0$) is present in
both specifications.  Three coefficients change sign
($\alpha_{\mathrm{Hispanic}}$, $\beta_{\mathrm{Black}}$,
$\beta_{\mathrm{Hispanic}}$), but the two intensive-margin sign changes
involve parameters whose Wald confidence intervals include zero, so the
sign itself is not statistically determined.  Overall, the comparison
confirms that the NSECE sampling design materially affects
point estimates---justifying the pseudo-posterior approach---while
leaving the central empirical findings intact.

\begin{table}[t]
\centering
\caption{Comparison of unweighted (M3b) and survey-weighted (M3b-W) fixed-effect
  posterior means. ``Shift'' is M3b-W minus M3b. ``Shift/SE'' normalizes the shift
  by the sandwich standard error, providing a scale-free measure of the weighting
  impact. ``Wald 95\% CI'' is the sandwich-corrected interval for M3b-W (\cref{thm:cholesky}).
  Asterisks indicate significance at the 5\% level. All covariates are standardized.}
\label{tab:m3b-comparison}
\smallskip
\small
\begin{adjustbox}{max width=\textwidth}
\begin{tabular}{@{}l rr rr l@{}}
\toprule
Parameter & M3b & M3b-W & Shift & Shift/SE & Wald 95\% CI \\
\midrule
\multicolumn{6}{@{}l}{\textit{Extensive margin ($\balpha$)}} \\[2pt]
Intercept            & $+0.558$ & $+0.764$ & $+0.206$ & $+4.64$ & $[0.677,\; 0.851]^{*}$ \\
Poverty              & $-0.220$ & $-0.324$ & $-0.103$ & $-1.97$ & $[-0.426,\; -0.221]^{*}$ \\
Urban                & $+0.264$ & $+0.442$ & $+0.178$ & $+5.19$ & $[0.375,\; 0.509]^{*}$ \\
Black                & $+0.248$ & $+0.478$ & $+0.230$ & $+5.20$ & $[0.392,\; 0.565]^{*}$ \\
Hispanic             & $-0.146$ & $+0.006$ & $+0.151$ & $+2.94$ & $[-0.095,\; 0.106]$ \\[4pt]
\multicolumn{6}{@{}l}{\textit{Intensive margin ($\bbeta$)}} \\[2pt]
Intercept            & $-0.096$ & $-0.242$ & $-0.146$ & $-9.65$ & $[-0.272,\; -0.213]^{*}$ \\
Poverty              & $+0.040$ & $+0.090$ & $+0.051$ & $+2.68$ & $[0.053,\; 0.127]^{*}$ \\
Urban                & $-0.055$ & $-0.047$ & $+0.008$ & $+0.39$ & $[-0.087,\; -0.007]^{*}$ \\
Black                & $+0.102$ & $-0.017$ & $-0.119$ & $-6.51$ & $[-0.052,\; 0.019]$ \\
Hispanic             & $+0.060$ & $-0.097$ & $-0.157$ & $-9.35$ & $[-0.130,\; -0.064]^{*}$ \\[4pt]
\multicolumn{6}{@{}l}{\textit{Overdispersion}} \\[2pt]
$\log\kappa$         & $+1.661$ & $+1.919$ & $+0.257$ & $+8.80$ & $[1.861,\; 1.976]^{*}$ \\
\bottomrule
\end{tabular}
\end{adjustbox}
\end{table}


\paragraph{The poverty reversal.}
The reversal is the central empirical finding. The extensive-margin
poverty coefficient is $\hat{\alpha}_{\mathrm{pov}} = -0.324$ (Wald 95\% CI:
$[-0.426, -0.221]$; $\DER = 2.41$): a one-standard-deviation increase in
community poverty (approximately 8.3 percentage points) reduces the log-odds
of serving infants by 0.324. The intensive-margin poverty coefficient is
$\hat{\beta}_{\mathrm{pov}} = +0.090$ (Wald 95\% CI:
$[+0.053, +0.127]$; $\DER = 1.61$): among centers that clear the hurdle,
the same poverty increase raises the logit IT share by 0.090. Both intervals
exclude zero, the signs are firmly opposed, and the posterior probability
that the reversal holds at the population average is
$\Pr(\alpha_{\mathrm{pov}} < 0 \text{ and } \beta_{\mathrm{pov}} > 0
\mid \text{data}) = 1.000$,
computed from the sandwich-corrected posterior draws
(the binding margin, $\beta_{\mathrm{pov}}$, lies $4.8$
sandwich standard deviations from zero; a bivariate normal
approximation gives a complementary probability of order~$10^{-6}$).

\paragraph{Naive versus Wald intervals: resolving prior domination.}
The contrast between the naive and Wald intervals in
\cref{tab:fixed-effects} provides a clear empirical demonstration of
the prior domination phenomenon. The naive
credible intervals for $\balpha$ and $\bbeta$ span approximately 3.7 units
on the logit scale---wide enough to include both positive and negative values
for every parameter, rendering significance testing impossible. The Wald
intervals, by contrast, are 0.06--0.21 units wide and permit clear
inferences.

This discrepancy arises because the naive posterior of the population-average
coefficients is dominated by the prior rather than the data. In the
hierarchical model, the state random effects $\bdelta_{k,s}$ absorb nearly
all data information about the state-specific coefficients
$\tilde{\alpha}_{k,s} = \alpha_k + \delta_{k,s}$
\citep[cf.][]{RabeHeskethSkrondal2006}; the population-average
$\alpha_k$ is identified only through the random effects distribution, with
an effective sample size of $S = 51$ (the number of states) rather than
$N = 6{,}785$ (the number of providers). Consequently, the marginal
posterior variance $[\bSigma_{\mathrm{MCMC}}]_{pp}$ remains close to the
prior variance rather than converging to the data-informed precision
$[\mathbf{H}_{\mathrm{obs}}^{-1}]_{pp}$. The hierarchical variance
ratio---the ratio of MCMC posterior variance to observed-information
variance---exceeds 600 for all fixed effects (see the computational remark
in~\cref{sec:survey}; a full diagnostic table appears in~\cref{smd:tab-der}).

The sandwich-corrected Wald intervals bypass this hierarchical inflation entirely.
The Wald standard error $\sqrt{[\mathbf{V}_{\mathrm{sand}}]_{pp}}$ is a
frequentist quantity derived from the observed information and the
cluster-robust outer product (\cref{eq:V-sand})
\citep{Binder1983, WilliamsSavitsky2021}: it measures the sampling
variability of $\hat{\btheta}$ across hypothetical replications of the
survey, not the width of the Bayesian posterior. For fixed effects in
hierarchical survey models, this frequentist precision is the appropriate
basis for population-level inference. The naive posterior intervals remain
useful for hyperparameters (where the sandwich is not applicable) and for
state random effects (where the prior appropriately regularizes small-sample
estimates).

\paragraph{Design effect ratios.}
The DER column in~\cref{tab:fixed-effects} reveals a substantively
interpretable pattern. The largest DER is 4.18 for
$\alpha_{\mathrm{Urban}}$: this coefficient is most severely inflated by
the survey design because the NSECE's stratification scheme
\citep{NSECEProjectTeam2022} is explicitly
linked to urbanicity---urban and rural areas are sampled at different rates.
The second-largest DER is 3.18 for $\beta_{\mathrm{Urban}}$, confirming
that the design effect is feature-specific rather than margin-specific. At
the other extreme, the smallest DER is 1.14 for $\beta_{\mathrm{Hisp}}$,
indicating that the survey design has minimal impact on this coefficient.
The overall mean DER of 2.11 is broadly consistent with the Kish $\DEFF$
of 3.76. \cref{fig:sandwich-impact} displays the naive and sandwich
intervals side by side for all 11 parameters, together with the DER
bar chart.

\begin{figure}[t]
  \centering
  \includegraphics[width=\textwidth]{F6_sandwich_correction.pdf}
  \caption{Impact of the sandwich variance correction on fixed-effect
    inference. \emph{Left panel:} naive MCMC 95\% credible intervals (gray)
    versus sandwich-corrected Wald 95\% confidence intervals (colored) for
    all 11 fixed effects. The naive intervals reflect prior width; the Wald
    intervals reflect data information adjusted for the survey design.
    \emph{Right panel:} design effect ratio ($\DER$) for each parameter. A
    $\DER$ of 1 indicates no survey design effect; the observed range is
    1.14--4.18 (mean 2.11).}\label{fig:sandwich-impact}
\end{figure}

\paragraph{Overdispersion.}
The estimated concentration parameter is $\hat{\kappa} = \exp(1.919)
= 6.81$ (Wald CI for $\log\kappa$: $[1.861, 1.976]$), confirming the
substantial overdispersion previewed in~\cref{sec:data}. With a typical
enrollment of $n_i \approx 50$, the variance inflation factor
$(n + \kappa)/(1 + \kappa) \approx 57/8 \approx 7$ is roughly sevenfold
the binomial baseline. Notably, $\log\kappa$ is the only parameter for
which the naive and Wald intervals are comparable in width (0.106 vs.\
0.115), because $\kappa$ has no corresponding state random effect---it is a
single global parameter informed directly by all $N_{\mathrm{pos}} = 4{,}392$
positive observations, so the standard Bernstein--von Mises approximation
holds and prior domination does not arise.

%% ---------------------------------------------------------------------------
%% 5.3 State-Specific Heterogeneity and Cross-Margin Independence
%% ---------------------------------------------------------------------------

\subsection{State-specific heterogeneity and cross-margin independence}\label{sec:statevar}

The population-average coefficients describe the national-level effects;
the state-varying coefficients reveal where and how those effects change
across jurisdictions.

\paragraph{Cross-margin scatter and independence.}
\cref{fig:cross-margin} plots the posterior means
$(\tilde{\alpha}_{\mathrm{pov},s},\;
\tilde{\beta}_{\mathrm{pov},s})$ across all 51 states. Under the
hierarchical M2 specification, 48 of 51 states (94\%) fall in the
upper-left quadrant ($\alpha < 0$, $\beta > 0$), compared with 23
(45\%) in the preliminary OLS analysis (\cref{sec:features})---a
dramatic increase driven by hierarchical shrinkage pulling
small-state estimates toward the negative population-average
extensive-margin coefficient. The posterior cross-margin correlation
for the poverty coefficient is
$\hat{\varrho}_{\mathrm{pov}}^{\mathrm{cross}} \approx 0.021$,
indistinguishable from zero, implying that the extensive and intensive
margins of the poverty effect operate through largely independent
mechanisms---a structural finding consistent with the block-diagonal
identification result of \cref{prop:sigma12-nonid}. Full state-specific coefficient estimates appear in
Supplementary Material~D (\cref{smd:states}); the reversal
probability map is provided in \cref{smd:reversal-map}.

\begin{figure}[t]
  \centering
  \includegraphics[width=0.85\textwidth]{F5_cross_margin_scatter.pdf}
  \caption{Cross-margin scatter plot of state-specific poverty coefficients
    from the unweighted M2 model (block-diagonal SVC).
    Each point represents the posterior mean
    $(\tilde{\alpha}_{\mathrm{pov},s},\;\tilde{\beta}_{\mathrm{pov},s})$ for
    one state.
    Small points without labels denote states with $N < 40$.  The four
    quadrants correspond to different poverty--enrollment patterns: the
    upper-left quadrant (shaded) contains 48 of 51 states (94\%),
    indicating a near-universal reversal pattern under hierarchical shrinkage.
    The dashed lines indicate the population-average coefficients
    $(\alpha_{\mathrm{pov}},\;\beta_{\mathrm{pov}})$. The near-zero
    cross-margin correlation ($\hat{\varrho} \approx 0.021$) implies that the
    two margins respond to poverty through largely independent
    mechanisms.}\label{fig:cross-margin}
\end{figure}

%% ---------------------------------------------------------------------------
%% 5.4 Marginal Effect Decomposition
%% ---------------------------------------------------------------------------

\subsection{Marginal effect decomposition}\label{sec:ame}

This subsection demonstrates the marginal effect decomposition---an
analytical framework enabled by the two-part structure---by translating
logit-scale coefficients into probability-scale quantities, building on
the log-scale decomposition in~\cref{prop:lae-lie}.  The decomposition
separates each covariate's total effect into extensive (access) and
intensive (share) components, a partition available only under the
two-part specification.  \Cref{tab:ame-decomp} reports the results.

\paragraph{From log to probability scale.}
\Cref{prop:lae-lie} decomposes the marginal effect on the
log-expected-enrollment scale, where the two channels are exactly
additive.  To translate to the probability scale, apply the product
rule to $E[Y_i/n_i] = q_i\,h_i$:
\begin{equation}\label{eq:ame-bridge}
  \frac{\partial\,(q_i\,h_i)}{\partial x_{i,k}}
  = \underbrace{h_i\,q_i(1-q_i)\,\tilde{\alpha}_{k,s}}_{\text{extensive}}
  \;+\;
  \underbrace{q_i\,h_i\,\varepsilon_{h,i}\,(1-\mu_i)\,
             \tilde{\beta}_{k,s}}_{\text{intensive}}
  \;=\; q_i\,h_i\bigl(\LAE_k + \LIE_k\bigr).
\end{equation}
The $\AME$ entries in~\cref{tab:ame-decomp} average each component
over providers and Cholesky-calibrated posterior
draws~$\btheta^{*(m)}$ (\cref{thm:cholesky}), mapping the log-scale
decomposition of~\cref{prop:lae-lie} onto the percentage-point scale.

\paragraph{Poverty.}
On the extensive margin, the $\AME$ of poverty is
$-0.029$ (95\% CI: $[-0.038, -0.020]$): a one-SD increase
in community poverty (${\approx}\,8.3$ percentage points)
reduces the probability of serving any infants by 2.9
percentage points.  On the intensive margin, $\AME = +0.015$
(95\% CI: $[+0.009, +0.021]$): among centers that do serve
infants, the same poverty increase raises the IT share by
1.5 percentage points.  The total $\AME$ is $-0.014$
(95\% CI: $[-0.027, -0.001]$), with the extensive margin
accounting for 66.0\% of the total absolute effect: the access
barrier dominates, but the intensity response offsets roughly
half of it.  This decomposition---unavailable under any
single-equation specification---illustrates the analytical
value of the two-part structure (\cref{prop:lae-lie}): the
hurdle identifies the access barrier as the dominant channel,
a finding that a single-equation model would collapse into a
muted net effect of $-1.4$~percentage points.

\paragraph{Other covariates.}
A notable pattern in~\cref{tab:ame-decomp} is that all four covariates
exhibit opposing extensive and intensive effects---the reversal pattern is
not unique to poverty. Urbanicity increases access ($+3.9$ pp) but decreases
intensity ($-0.8$ pp), with the extensive margin accounting for 83.6\% of
the total. The Black population share increases access ($+4.3$ pp) with a
negligible intensity effect (Ext.\ Share $= 93.9\%$). Hispanic composition
stands out: the extensive $\AME$ is effectively zero ($+0.001$), so the
entire effect operates through the intensive margin, where higher Hispanic
composition is associated with a lower IT share ($-1.6$ pp). This
intensive-margin-dominated pattern contrasts sharply with the
extensive-margin-dominated patterns for the other covariates. In every case,
a single-equation model would report only the net total effect, concealing
the offsetting channels---precisely the decomposition that the hurdle
framework and the $\AME$ formulas in \cref{prop:lae-lie} are designed to
reveal.

\begin{table}[t]
\centering
\caption{Average marginal effect ($\AME$) decomposition from model M3b.
  Effects are on the probability scale (change per one-SD increase in the
  covariate). ``Ext.\ Share'' is the percentage of the total absolute
  effect attributable to the extensive (access) margin, computed as
  $|\text{Extensive}|/(|\text{Extensive}| + |\text{Intensive}|)$.
  All four covariates
  exhibit opposing extensive and intensive effects.  The 95\% intervals
  are 2.5th and 97.5th posterior percentiles of each $\AME$ component
  evaluated at the Cholesky-calibrated
  draws~$\btheta^{*(m)}$~(\cref{thm:cholesky}).}\label{tab:ame-decomp}
\smallskip
\begin{tabular}{@{}lcccc@{}}
\toprule
Covariate & Extensive & Intensive & Total & Ext.\ Share \\
\midrule
Poverty & $-0.029$ & $+0.015$ & $-0.014$ & 66.0\% \\
 & $[-0.038,\,-0.020]$ & $[+0.009,\,+0.021]$ & $[-0.027,\,-0.001]$ & \\[3pt]
Urban & $+0.039$ & $-0.008$ & $+0.032$ & 83.6\% \\
 & $[+0.034,\,+0.045]$ & $[-0.014,\,-0.001]$ & $[+0.022,\,+0.041]$ & \\[3pt]
Black & $+0.043$ & $-0.003$ & $+0.040$ & 93.9\% \\
 & $[+0.035,\,+0.051]$ & $[-0.009,\,+0.003]$ & $[+0.029,\,+0.052]$ & \\[3pt]
Hispanic & $+0.001$ & $-0.016$ & $-0.016$ & 2.8\% \\
 & $[-0.009,\,+0.010]$ & $[-0.021,\,-0.010]$ & $[-0.027,\,-0.004]$ & \\
\bottomrule
\end{tabular}
\end{table}

%% ---------------------------------------------------------------------------
%% 5.5 Policy Moderation
%% ---------------------------------------------------------------------------

\subsection{Policy moderation}\label{sec:policyresults}

The reversal rate and the geographic clustering documented in
\cref{sec:statevar} raise a natural question: can observable CCDF policies
explain why the reversal is strong in some states and absent in
others?  The $\bGamma_k$ matrices in M3b decompose the state-level
variation in covariate effects into policy-explained and residual
components (\cref{prop:poverty-expansion}).  The full $40$-element
$\bGamma$ matrices appear in Supplementary Material~D
(\cref{smd:gamma}); here we summarize three key patterns.

First, tiered reimbursement attenuates the extensive-margin poverty
barrier ($\hat{\gamma}_{1,\mathrm{pov},\mathrm{TR}} = +0.495$,
$\Pr(>0) = 0.962$) and raises baseline intensive-margin enrollment
($\hat{\gamma}_{2,\mathrm{int},\mathrm{TR}} = +0.271$,
$\Pr(>0) = 0.985$).  Second, IT rate add-ons raise baseline
participation ($\hat{\gamma}_{1,\mathrm{int},\mathrm{IT}} = +0.728$,
$\Pr(>0) = 0.996$) but paradoxically \emph{strengthen} the poverty
gradient on the extensive margin
($\hat{\gamma}_{1,\mathrm{pov},\mathrm{IT}} = -0.488$,
$\Pr(<0) = 0.982$), suggesting that the add-on primarily benefits
moderate-poverty areas.  Third, IT add-ons are most effective at closing
demographic gaps in infant enrollment, particularly for centers in Black
communities on the intensive margin
($\hat{\gamma}_{2,\mathrm{Black},\mathrm{IT}} = +0.256$,
$\Pr(>0) = 1.000$).

The three observable policy variables explain between 4\% and 16\% of
the between-state variation in poverty coefficients.  These associations
are cross-sectional and do not support causal claims; nonetheless, the
$\bGamma$ decomposition illustrates how the hierarchical structure in
\cref{sec:model} partitions state heterogeneity into policy-attributable
and residual components.
